\section{The endomorphism algebra and homological consequences}
\label{sec:consequences}

The symmetric pair constructed in Theorem~\ref{thm:main} has a
self-orthogonal generator $A\oplus M$. The classical relation between
such generators and dominant dimension turns the module counterexample
into an algebra whose entire regular injective resolution is projective.
We first isolate this step, then extract the modules in
Corollaries~\ref{cor:homological} and~\ref{cor:finitistic}.

\begin{proposition}\label{prop:endomorphism}
Let $B$ be a finite-dimensional symmetric algebra over a field and let
$N$ be a finite-dimensional nonprojective left $B$-module satisfying
$\Ext_B^i(N,N)=0$ for all $i>0$. Put
\[
 G=B\oplus N,\qquad \Gamma=\End_B(G)^{\mathrm{op}},
\]
where endomorphisms compose in the ordinary order. Then
\[
 \domdim({}_\Gamma\Gamma)=\infty,
 \qquad \idim_\Gamma\Gamma=\infty.
\]
In particular, $\Gamma$ is not self-injective.
\end{proposition}

\begin{proof}
The module $G$ is a generator-cogenerator: it contains the regular
module, and symmetry identifies $B$ with its dual, an injective
cogenerator. Moreover
\[
 \Ext_B^i(G,G)=0\qquad(i>0).
\]
The $N,N$ summand vanishes by hypothesis. The summands with first
argument $B$ vanish by projectivity, and the remaining summand with
second argument $B$ vanishes by injectivity.

Mueller's criterion~\cite[Lemma~3]{Mueller1968} says that, for a
generator-cogenerator $G$, vanishing of these groups for
$1\leq i\leq n$ is equivalent to dominant dimension at least $n+2$
for its endomorphism algebra. His Theorem~4 identifies the left and
right dominant dimensions, so the opposite convention in the
definition of $\Gamma$ makes no difference to this assertion.
Applying the criterion for every $n$ proves
$\domdim({}_\Gamma\Gamma)=\infty$.

We verify non-self-injectivity directly. The functor
\[
 H=\Hom_B(G,-):\add G\longrightarrow \Gamma\text{-}\mathrm{proj}
\]
is fully faithful and sends $G$ to the left regular module $\Gamma$.
Explicitly, if $e\in\End_B(G)$ and $h:G\to V$, then
$e^{\mathrm{op}}h=h\circ e$; full faithfulness follows first on
finite sums of $G$ and then on their summands.
Suppose that the left regular module $\Gamma$ were injective. Its direct
summand $H(N)$ would be injective as well. Choose a finite injective envelope
$N\hookrightarrow J$ over $B$. Since $B$ is symmetric,
$J\in\add B\subseteq\add G$. Left exactness of $H$ gives a
monomorphism $H(N)\hookrightarrow H(J)$, which splits by injectivity
of $H(N)$. Full faithfulness gives a retraction $J\to N$ of the
original inclusion. Thus $N$ is injective, hence projective over the
symmetric algebra $B$, a contradiction.

Finally, every term in the minimal injective resolution of $\Gamma$
is projective by infinite dominant dimension. If its injective
dimension were finite, this resolution would end in a finite exact
sequence
\[
 0\longrightarrow\Gamma\longrightarrow I^0\longrightarrow\cdots
 \longrightarrow I^d\longrightarrow0
\]
of projective-injective terms. For $d>0$, the last surjection splits
because $I^d$ is projective. Its kernel is then projective, and
repeating backwards shows that the first inclusion splits. The case
$d=0$ is immediate. In either case $\Gamma$ is a summand of $I^0$
and is injective, contrary to the preceding paragraph. Hence
$\idim_\Gamma\Gamma=\infty$.
\end{proof}

\begin{proof}[Proof of Corollary~\ref{cor:homological}]
Fix a field extension $K/k$. The symmetric structure survives
extension of scalars: since $A$ is finite-dimensional, the natural
bimodule identification
\[
 K\otimes_k DA\simeq D_K(K\otimes_k A),
 \qquad D_K=\Hom_K(-,K),
\]
turns $A\simeq DA$ into a symmetrizing isomorphism for $A_K$.

We also record why the module hypotheses survive. For finite
$A$-modules $U,V$, resolve $U$ by finite projective modules.
Finite projectivity gives, in every degree,
\[
 K\otimes_k\Hom_A(P_i,V)
 \simeq
 \Hom_{A_K}(K\otimes_kP_i,K\otimes_kV).
\]
Flatness commutes with kernels and images, so
\begin{equation}\label{c:basechange}
 K\otimes_k\Ext_A^i(U,V)
 \simeq \Ext_{A_K}^i(K\otimes_kU,K\otimes_kV)
 \qquad(i\geq0).
\end{equation}
Thus $M_K$ is self-orthogonal. Choose a finite projective surjection
$P\twoheadrightarrow M$ with kernel $L$. Its extension class in
$\Ext_A^1(M,L)$ is nonzero, since $M$ is not projective.
It remains nonzero under \eqref{c:basechange}, because extension of
fields is faithful. The extended surjection therefore does not split,
and $M_K$ is nonprojective. The degree-zero case of
\eqref{c:basechange} is compatible with composition and also identifies
\[
 \Gamma_K\simeq K\otimes_k\End_A(A\oplus M)^{\mathrm{op}}.
\]
These are the same finite Ext and nonprojectivity base-change
arguments used in the companion~\cite[Proposition~9.2]{OpenAIAR2026}.

Apply Proposition~\ref{prop:endomorphism} to $A_K,M_K$, and abbreviate
$\Gamma_K$ to $\Gamma$. Write
\[
 0\longrightarrow\Gamma\longrightarrow I^0
 \longrightarrow I^1\longrightarrow\cdots
\]
for its minimal injective resolution. Every $I^j$ is
projective-injective and $\idim_\Gamma\Gamma=\infty$.
This proves the first assertion and refutes the classical Nakayama
conjecture. It also refutes the Auslander--Gorenstein conjecture:
$\pd_\Gamma I^j=0\leq j$ for every $j\geq0$, whereas the left
self-injective dimension is already infinite.

There is a nonprojective indecomposable injective left $\Gamma$-module
$J$. Indeed, finite-dimensional algebras have the same finite number
of isomorphism classes of indecomposable projectives and indecomposable
injectives, both indexed by the simple modules. If every indecomposable
injective were projective, these classes would exhaust the indecomposable
projectives, making the regular module injective. This contradicts
Proposition~\ref{prop:endomorphism}. Since every $I^j$ is projective,
none contains $J$ as a direct summand, giving the omitted-injective
failure of the generalized Nakayama conjecture.

Let $S=\soc J$. The socle of an indecomposable injective over a
finite-dimensional algebra is simple, and $J$ is its injective
envelope. If there were a nonzero map $S\to I^j$, it would be
injective. Injectivity of $I^j$ would extend it to a map $J\to I^j$.
The extension is injective because $S$ is essential in $J$, and it
splits because $J$ is injective. This would make $J$ a summand of
$I^j$, a contradiction. Hence $\Hom_\Gamma(S,I^j)=0$ for every
$j\geq0$. Applying $\Hom_\Gamma(S,-)$ to the resolution gives
\[
 \Ext_\Gamma^i(S,\Gamma)=0\qquad(i\geq0).
\]
The nonzero module $S$ therefore refutes the strong Nakayama conjecture
and also gives the simple-Ext formulation of generalized Nakayama
failure. In particular, the degree-zero vanishing is part of the
conclusion, not only the positive-degree vanishings.

Finally let $W$ be the direct sum of one representative from each
isomorphism class of indecomposable projective-injective left
$\Gamma$-modules. There are finitely many such classes, so $W$ is
finite, projective, and injective. In particular
$\Ext_\Gamma^i(W,W)=0$ for $i>0$, and every $I^j$ belongs to
$\add W$. The displayed injective resolution is therefore an
infinite $\add W$-coresolution of $\Gamma$. Applying
$\Hom_\Gamma(-,W)$ preserves its exactness because $W$ is injective.
Thus $W$ is Wakamatsu tilting and has projective dimension zero.

If $W$ were tilting, there would be a finite exact coresolution
\[
 0\longrightarrow\Gamma\longrightarrow W^0\longrightarrow\cdots
 \longrightarrow W^r\longrightarrow0,\qquad W^j\in\add W.
\]
Every $W^j$ is projective-injective. The backward splitting argument
in Proposition~\ref{prop:endomorphism} would make $\Gamma$ injective,
a contradiction. Hence $W$ is not tilting, which refutes the Wakamatsu
tilting conjecture.
\end{proof}

The proof works separately over each extension $K/k$. It does not
require a minimal injective resolution to remain minimal under scalar
extension, or an indecomposable injective or simple module to remain
indecomposable or simple. The construction of $\Gamma_K$ is compatible
with scalar extension, and the required $J,S,W$ are chosen over $K$.

The successive cokernels of the same injective resolution give the
finitistic-dimension consequence.

\begin{proof}[Proof of Corollary~\ref{cor:finitistic}]
Fix a field extension $K/k$ and abbreviate $\Gamma_K$ to $\Gamma$.
Use the minimal injective resolution of the left regular module over
this field,
\[
 0\longrightarrow {}_\Gamma\Gamma\longrightarrow I^0
 \longrightarrow I^1\longrightarrow\cdots.
\]
Every $I^n$ is finite-dimensional and projective-injective by
Corollary~\ref{cor:homological}. Set $C_0={}_\Gamma\Gamma$ and
$C_n=\im(I^{n-1}\to I^n)$ for $n\geq1$. Exactness gives
\[
 0\longrightarrow C_n\longrightarrow I^n\longrightarrow C_{n+1}
 \longrightarrow0\qquad(n\geq0).
\]
Thus the $C_{n+1}$ are successive cokernels, and all $C_n$ are finite.

Since $C_0$ and $I^0$ are projective, the first sequence gives
$\pd_\Gamma C_1\leq1$. If $C_1$ were projective, its surjection from
$I^0$ would split, making $\Gamma$ a direct summand of the injective
module $I^0$. This would make $\Gamma$ injective, contrary to
$\idim_\Gamma\Gamma=\infty$. Hence $\pd_\Gamma C_1=1$.

Suppose $n\geq1$ and $\pd_\Gamma C_n=n$. Splicing a length-$n$
projective resolution of $C_n$ with the displayed short exact sequence
gives $\pd_\Gamma C_{n+1}\leq n+1$. For every left module $V$ and
$j\geq1$, projectivity of $I^n$ and the long exact Ext sequence give
\[
 \Ext_\Gamma^j(C_n,V)\simeq\Ext_\Gamma^{j+1}(C_{n+1},V).
\]
There is a left module $V$ with $\Ext_\Gamma^n(C_n,V)\neq0$, since
otherwise the Ext criterion for projective dimension would give
$\pd_\Gamma C_n\leq n-1$. Thus
$\Ext_\Gamma^{n+1}(C_{n+1},V)\neq0$, proving the reverse bound and
$\pd_\Gamma C_{n+1}=n+1$. Induction gives the asserted exact dimensions.
The finite left modules $C_n$ therefore make the little left finitistic
dimension infinite. The big left finitistic dimension is at least the
little one, so it is infinite as well.
\end{proof}

The remaining applications use the established failures of the homological
conjectures together with results on tilting modules and stratifying ideals.

A finite left module $P$ over a finite-dimensional algebra $\Lambda$
is \emph{almost complete tilting} if it is a direct summand of a tilting
module and has one fewer isomorphism class of indecomposable summands
than there are isomorphism classes of simple left $\Lambda$-modules.
Here tilting has the ordinary finite-projective-dimension meaning
recalled in Section~\ref{sec:introduction}. An \emph{indecomposable
complement} is a finite indecomposable module $Y$ such that $P\oplus Y$
is tilting.

\begin{corollary}[Infinitely many tilting complements]\label{cor:tilting-complements}
For every field extension $K/k$, let $\Gamma_K$ be the algebra in
Corollary~\ref{cor:homological}. There is a finite-dimensional projective
almost complete tilting left $\Gamma_K^{\mathrm{op}}$-module $P_K$ with
infinitely many pairwise nonisomorphic indecomposable complements.
\end{corollary}

\begin{proof}
Fix $K$ and abbreviate $\Gamma_K$ to $\Gamma$. By
Corollary~\ref{cor:homological}, $\Gamma$ fails the generalized Nakayama
conjecture. Buan--Solberg~\cite[Theorems~3.1 and~3.4]{BuanSolberg1998}
therefore give an injective almost complete cotilting left $\Gamma$-module
$U$ with infinitely many pairwise nonisomorphic indecomposable complements
$X$. The exact duality $D_K=\Hom_K(-,K)$ sends finite left
$\Gamma$-modules to finite left $\Gamma^{\mathrm{op}}$-modules.
It sends finite injective dimension to finite projective dimension,
preserves self-orthogonality, and dualizes the finite cotilting resolution
of the injective cogenerator into a finite tilting coresolution of the
regular module. Thus $P_K=D_KU$ is projective and almost complete tilting,
and $D_KX$ are its complements. Duality preserves indecomposability and
distinct isomorphism classes, so there are infinitely many such
complements. The witnesses are chosen separately over each extension
$K/k$.
\end{proof}

The final application concerns decompositions of derived categories.
For an idempotent $e$ in a finite-dimensional algebra $R$, the ideal
$ReR$ is \emph{stratifying} if multiplication induces an isomorphism
\[
 Re\otimes_{eRe}eR\xrightarrow{\sim}ReR
\]
and $\Tor_i^{eRe}(Re,eR)=0$ for every $i>0$.
It then induces the standard recollement relating the unbounded
derived categories of $R/ReR$, $R$, and $eRe$; see
\cite[Introduction]{ChenFangXi2024}.
An \emph{idempotent stratification of length $n$}, in the convention of
\cite[Definition~4.7]{ChenFangXi2024}, consists of nonzero mutually
orthogonal idempotents $e_0,\ldots,e_n$ with sum $1$ such that, on putting
$f_i=e_0+\cdots+e_i$, the ideal generated by $f_i$ in
$f_{i+1}Rf_{i+1}$ is stratifying and $e_{i+1}\notin Rf_iR$ for
$0\leq i<n$. Thus its length counts the successive nontrivial
recollements obtained by passing through these corner algebras.

\begin{corollary}[Idempotent stratifications]\label{cor:stratifying}
For every field extension $K/k$ and every integer $L\geq1$, there is an
auxiliary finite-dimensional indecomposable symmetric $K$-algebra $T_L$
admitting an idempotent stratification of length at least $L$ in the sense
of Chen, Fang and Xi. In particular, $T_L$ has a nonzero proper
stratifying ideal, and $\mathcal D(T_L\text{-}\mathrm{Mod})$ admits a
chain of at least $L$ nontrivial standard recollements through successive
corner algebras. Here $\mathrm{Mod}$ denotes all left modules and
$\mathcal D$ is the unbounded derived category.
\end{corollary}

\begin{proof}
The base-change argument above gives a symmetric $K$-algebra $A_K$ and a
nonprojective self-orthogonal module $M_K$. Choose a nonprojective
indecomposable summand $N$ in one block $B$ of $A_K$. The block $B$ is
indecomposable and symmetric, and $\Ext_B^i(N,N)=0$ for every $i>0$,
since these groups are direct summands of $\Ext_{A_K}^i(M_K,M_K)$.

Chen, Fang and Xi~\cite[proof of Theorem~5.13]{ChenFangXi2024} construct
from such a pair, over any field, reduced mirror-reflective algebras $T_n$
that are indecomposable and symmetric. Their Corollary~5.12(1), with
Definition~4.7, gives an idempotent stratification of length at least
$n$. Theorem~1.2(2) supplies the
stated successive nontrivial standard recollements of unbounded derived
categories. Taking $n=L$ proves the claim.
\end{proof}
